\chapter{Introduction}
\label{sec1:introduction}

The rapid global emergence of antimicrobial resistance (AMR) represents one of the most critical threats that modern 
medicine is facing. As conventional antibiotics lose effiency against multi drug resistant bacterial pathogens, alternative 
therapeutic strategies have gained widespread attention. Among these is phage therapy, it is the targeted clinical application of 
bacteriophages (viruses that specifically infect and destroy bacteria) and has emerged as a highly promising solution 
\cite{dejonge2019molecular}.

Bacteriophages are the most abundant biological entities in the global biosphere, with an estimated population exceeding 
$10^{31}$ particles \cite{dejonge2019molecular}. One of the main determinants of a phage's therapeutic viability is its host range 
(the specific spectrum of bacterial strains it can successfully infect). Molecular recognition is governed by highly specialized 
interactions between viral structural proteins (such as tail fibers, lysins etc.) and bacterial cell-wall receptors, therefore 
phages are extremely host specific. While this specificity is a major clinical asset that allows therapeutic phages to eradicate 
targeted pathogens without harming the patient, it also introduces a diagnostic challenge, that is matching a bacterial infection 
with an effective phage.

Traditional laboratory screening methods, such as plaque assays, are labor intensive, costly, and scale poorly against 
the millions of uncharacterized viral genomes generated by modern high throughput metagenomic sequencing. This neccessiates 
the development of fast and scalable computational tools, that are capable of predicting host ranges directly from viral genomic 
and protein-level representations \cite{coclet2021global}.

However, state of the art computational prediction tools suffer from a fundamental limitation. They produce point 
predictions or uncalibrated scalar scores in $[0,1]$ without rigorous, finite sample uncertainty quantification. In high stakes 
medical applications, point predictions offer no guarantee of reliability. A confidence score of $0.85$ reflects a relative 
ranking not a calibrated conditional probability $\mathbb{P}(Y = y \mid X = x)$. When a false negative means missing a life-saving 
treatment, such predictions are clinically insufficient

\newpage
The main objective of this work is to resolve this limitation by contructing an uncertainty aware statistical framework 
for phage-host prediction using conformal prediction \cite{vovk2005algorithmic, angelopoulos2023conformal},which is a 
distribution free method that transforms point scoring models into set valued predictions $C(X) \subseteq \mathcal{H}$. 
More importantly, these prediction sets are mathematically guaranteed to contain the true host label with a user specified 
confidence level $1-\alpha$ (e.g., $90\%$), which holds strictly in finite samples under just the assumption of data exchangeability.

In this work we build on the Conformal Set Aggregation (CSA) framework \cite{ochoarivera2024conformal}. We address 
three major challenges that arise when applying CSA to classification datasets:

\begin{enumerate}[label=(\roman*)]
    \item \textbf{Non-convex geometry in classification space:} We demonstrate why traditional convex envelopes, developed 
           primarily for continuous regression spaces, are inefficient in classification settings. Here class conditional score 
           vectors occupy disjoint, bounded clusters in $[0,1]^K$. So, we propose non-convex (Radial and Strip) and Collapsed 1D 
           envelope geometries that adapt to cluster topologies, and eliminate unpopulated empty space, which helps us reduce 
           prediction set sizes.
    
    \item \textbf{Closed form analytical calibration ($\tau$):} We derive exact, closed form analytical expressions for the 
           per point inclusion scaling function $\tau(\mathbf{s})$ across all envelope geometries. This eliminates the 
           computationally expensive numerical binary search required by standard CSA, which reduces the calibration time 
           complexity from $O(|S_2| \cdot N_{\text{iter}})$ to a single forward pass and $O(|S_2| \log |S_2|)$ sort.
    
    \item \textbf{Class Conditional Validity:} Global calibration can yield under coverage for complex host classes and over 
           coverage for simpler ones. To prevent this we incorporate per class quantile calibration ($\hat{t}_c$). This guarantees 
           exact $1-\alpha$ coverage conditionally for every host label individually rather than just on average.
    
    % \item \textbf{Structural Missingness \& Homology Contamination:} We formalize structural missingness in protein score vectors 
    % as a first-class feature rather than imputing artificial values. Furthermore, we implement protein-level cross-validation 
    % masking to eliminate optimistic score bias caused by homologous protein cluster overlap across genome-level splits.
\end{enumerate}

We evaluated our methodology systematically on two tasks of increasing complexity. First, 
a controlled binary Gram-type classification task ($N=2$ classes) to validate the geometric envelope constructions, and then on a 
multi class host-level prediction task ($N=54$ bacterial genera) using real-world protein language model embeddings.


The remainder of this report is organized as follows: Chapter~\ref{sec2: background} reviews the biological background 
of bacteriophages, computational state-of-the-art methods, and the mathematical foundations of split conformal prediction and CSA. 
Chapter~\ref{sec3: methodology} formalizes the prediction tasks, develops the non-convex envelopes, the nonconformity 
transformation, and the analytical derivation of $\tau(\mathbf{s})$. Chapter~\ref{sec4: data} details the dataset architecture, 
protein to phage NaN aware mean pooling operators, and cross split protein cluster contamination filtering. Chapter~
\ref{sec5: experiments} presents synthetic experiments and empirical evaluation results for both gram-type and host-level 
classification tasks. Finally, chapter~\ref{sec6: conclusion} summarises our findings, discusses the biological limitations, and 
outlines directions for future research.
%%%%%%%%%%%%%%%%%%%%%%%%%%%%%%%%%%%%%%%%%%%%%%%%%%%%%%%%%%%%%%%%%%%%%%%%%%%%%%%%%%%%%


